\bookchapter{제43장: 대화 엔지니어링 — AI와 효율적으로 협업하는 대화 기술}

\chapterintro{대화는 AI 코딩의 핵심 상호작용 방식이에요. 하지만 주의하세요: “대화”는 “잡담”이 아니에요. — 효율적인 AI 코딩 대화는 구조를 갖춘 커뮤니케이션 기술이고, 우리는 이를 “대화 엔지니어링”이라 불러요. 다루는 주제는 이것입니다: 최소한의 대화 횟수로 최고 품질의 코드 결과물을 얻는 법. 쉽게 말해 이번 장은 AI와 “말을 명확하게 주고받는” 방법을 가르쳐 주는 장이에요.}

\booksection{프롬프트 계층화 전략}

AI와의 대화는 세 계층으로 나눌 수 있고, 각 계층은 서로 다른 기능을 맡아요. 이 세 계층을 이해하면 블록 쌓듯 효율적인 프롬프트를 조합할 수 있어요.

\begin{center}
\begin{tabularx}{\linewidth}{|X|X|X|X|}
\hline
\textbf{계층} & \textbf{기능} & \textbf{예시} & \textbf{빈도} \tabularnewline
\hline
시스템 계층 & AI의 역할과 전역 제약 설정 & “너는 React 전문가이고, 코드는 반드시 TypeScript로 쓴다” & 프로젝트 시작 시 한 번 설정 \tabularnewline
작업 계층 & 지금 수행할 구체적 작업 기술 & “사용자명과 비밀번호 필드가 있는 폼 컴포넌트를 구현해 줘” & 각 작업 시작 시 제시 \tabularnewline
피드백 계층 & AI 결과물 평가와 수정 & “비밀번호 필드에 표시/숨기기 전환이 빠졌어. 추가해 줘” & 결과물 품질에 따라 제공 \tabularnewline
\hline
\end{tabularx}
\end{center}

\begin{center}\small\color{gray}대화 프롬프트의 세 계층\end{center}

시스템 계층 프롬프트는 보통 프로젝트 기억 파일(CLAUDE.md 등)에 적기 때문에 매 대화마다 반복할 필요가 없어요. 작업 계층은 대화에서 가장 자주 쓰는 층이에요. — “무엇을 할지”, “어떤 기술을 쓸지”, “어떤 제약이 있는지” 세 요소를 담아야 하죠. 피드백 계층은 실력이 가장 드러나는 부분이에요. — 좋은 피드백은 AI가 정확하게 수정하게 만들지만, 나쁜 피드백은 AI를 다시 추측하게 만들 뿐이에요.

\begin{tcbinfo}{꿀팁: 피드백 3단 콤보}
실용적인 피드백 템플릿은 “문제 지적 + 기대 설명 + 맥락 제공”이에요. 예를 들어: “25번 줄의 오류 처리가 console.log만 쓰고 있어(문제 지적). 프로덕션 환경에서는 통합 오류 핸들러로 사용자 친화적 메시지를 돌려줘야 해(기대 설명). 프로젝트에 handleError 함수가 있으니 그걸 쓸 수 있어(맥락 제공)”. 이런 피드백은 “오류 처리 좀 고쳐 줘”보다 열 배는 효과적이에요.

\end{tcbinfo}

\booksection{디버깅 대화: AI가 빨리 원인을 찾도록 오류 설명하기}

디버깅은 AI 코딩에서 가장 빈번한 상황 중 하나예요. 버그를 만났을 때 문제를 AI에 어떻게 설명하느냐가, AI가 빨리 원인을 찾고 고쳐 줄 수 있는지를 결정해요. 다음은 디버깅 대화의 네 가지 요점이에요:

\begin{enumerate}
  \item AI에 당신의 전달이 아니라 원본 정보를 “보여” 주세요. 오류 로그, 예외 스택, 관련 코드를 @ 참조로 AI에 바로 넘기고, “뭔가 에러가 나는 것 같아요” 식으로 설명하지 마세요. AI는 당신의 이해를 보는 것보다 원본 정보를 보는 편이 훨씬 정확해요.
  \item “기대 동작”과 “실제 동작”의 차이를 설명하세요. “버튼을 누르면 홈으로 이동할 거라고 기대했는데(기대), 실제로는 눌러도 페이지가 반응하지 않아요(실제)” — 이런 설명이 “버튼이 고장 났어요”보다 훨씬 유용해요.
  \item 재현 단계를 제공하세요. “먼저 사용자명을 입력하고, 로그인 버튼을 누르면 이 오류가 나타나요” — 재현 단계가 있으면 AI가 당신의 조작 경로를 따라가 볼 수 있어 문제를 더 쉽게 찾을 수 있어요.
  \item 이미 무엇을 시도했는지 분명히 말하세요. “네트워크 요청을 확인했고 200이 돌아왔으며 데이터도 정상이에요” — 이렇게 하면 AI가 이미 검증된 방향을 배제하고 같은 걸 반복해 조사하며 시간을 낭비하지 않아요.
\end{enumerate}

나쁜 디버깅 설명

\begin{lstlisting}
버튼을 눌러도 반응이 없어요. 좀 봐 주세요
\end{lstlisting}

좋은 디버깅 설명

\begin{lstlisting}
로그인 버튼을 클릭한 후 페이지가 반응하지 않음. 기대 동작: 클릭 후 /home 페이지로 이동 실제 동작: 클릭해도 아무 반응이 없고 콘솔에도 오류가 없음 재현 단계: 1. 로그인 페이지에서 사용자명 admin, 비밀번호 123456 입력 2. “로그인” 버튼 클릭 3. 페이지가 반응 없음 확인 완료: - 네트워크 요청 정상, 200 반환 - onClick 이벤트 바인딩됨 - LoginButton 컴포넌트 코드 확인함 관련 파일: @src/components/LoginButton.tsx @src/pages/Login.tsx
\end{lstlisting}

\booksection{반복 전략: 작은 단계 빠르게 vs 큰 단계 리팩터링}

AI 코딩의 반복에는 서로 다른 상황에 맞는 두 가지 기본 전략이 있어요. 둘의 차이를 이해하면 상황에 따라 올바른 선택을 할 수 있어요.

“작은 단계 빠르게” 전략은 기능 개발과 버그 수정에 어울려요. 핵심은: 매번 AI에 작은 수정 하나만 시키고, 검증을 통과한 뒤 다음으로 넘어가는 것. 장점은 매 단계가 검증 가능하고 되돌릴 수 있어 리스크를 통제할 수 있다는 점, 단점은 속도가 느리고 대화 왕래가 잦다는 점이에요. 어울리는 경우: 신규 기능 개발, 최적 방안이 확실치 않아 탐색할 때, AI 결과물이 좀 불안할 때의 신중한 반복이에요.

“큰 단계 리팩터링” 전략은 목표가 분명한 대규모 수정에 어울려요. 핵심은: 완전한 목표를 한 번에 설명하고 AI가 수정 전체를 스스로 계획·실행하게 하는 것. 장점은 속도가 빠르고 AI가 전체 그림을 보고 더 나은 아키텍처 결정을 내릴 수 있다는 점, 단점은 리스크가 높아 AI가 잘못 이해하면 수정 비용이 크다는 점이에요. 어울리는 경우: 명확한 아키텍처 업그레이드, 대규모 코드 이전, 충분한 테스트가 뒷받침되는 리팩터링이에요.

\begin{center}
\begin{tabularx}{\linewidth}{|X|X|X|}
\hline
\textbf{구분} & \textbf{작은 단계 빠르게} & \textbf{큰 단계 리팩터링} \tabularnewline
\hline
1회 수정량 & 기능 하나 또는 버그 하나 & 모듈 하나 또는 아키텍처 계층 하나 \tabularnewline
검증 빈도 & 수정할 때마다 즉시 검증 & 전체 수정 완료 후 일괄 검증 \tabularnewline
대화 횟수 & 많음(기능당 5\textasciitilde{}10회) & 적음(모듈당 1\textasciitilde{}3회) \tabularnewline
리스크 & 낮음, 매 단계 되돌리기 가능 & 높음, 오해 시 수정 비용 큼 \tabularnewline
적합한 상황 & 신규 기능, 탐색적 개발 & 목표가 분명한 리팩터링·이전(마이그레이션) \tabularnewline
전제 조건 & 특별한 요구 없음 & 테스트 커버리지 필요 \tabularnewline
\hline
\end{tabularx}
\end{center}

\begin{center}\small\color{gray}두 가지 반복 전략 비교\end{center}

실전에서는 두 전략을 조합해 쓰는 경우가 많아요. 예컨대 모듈 하나를 리팩터링할 때 먼저 “큰 단계 리팩터링”으로 AI에게 전체 방안을 제시받고, 방안을 검토한 뒤 “작은 단계 빠르게”로 하나씩 실행하는 식이에요. 이렇게 하면 전체 시야를 지키면서 실행 리스크도 통제할 수 있어요.

\booksection{여러 차례 대화의 컨텍스트 관리}

대화가 거듭될수록 대화 컨텍스트를 능동적으로 관리해 “AI가 초반 정보를 잊어버리는” 문제를 피해야 해요. 다음은 실용적인 컨텍스트 관리 기법 몇 가지예요.

첫째, 핵심 정보를 “닻으로 고정”하세요. 대화를 시작할 때 핵심 제약을 분명히 적고, 이후 대화에서 적절히 다시 언급해요. 예컨대 “기억해 줘, 우리 데이터는 localStorage에 저장하고 데이터베이스는 쓸 수 없어” — 관련 수정 요청에서 이 제약을 다시 상기시켜 AI가 잊지 않게 하세요.

둘째, “요약-이어가기” 기법을 잘 쓰세요. 대화가 아주 길어지면 AI에게 “지금까지 진행 상황과 남은 할 일을 요약해 줘”라고 하고, 새 대화를 열어 그 요약을 출발점으로 삼아요. 이렇게 하면 핵심 컨텍스트를 지키면서 긴 대화의 성능 저하도 피할 수 있어요.

셋째, “한 대화에 한 작업”. 가장 단순하면서도 가장 효과적인 컨텍스트 관리 전략이에요. 대화 하나당 작업 하나, 기능 하나만 처리하고 끝나면 새 대화에서 다음을 처리하세요. 대화 길이가 자연스럽게 제한되어 컨텍스트가 깔끔하게 유지돼요.

넷째, 도구의 컨텍스트 관리 기능을 활용하세요. 대부분의 AI 코딩 도구에는 “컨텍스트 지우기” 또는 “새 세션” 버튼이 있어요. — AI의 답변 품질이 떨어진다고 느껴지면 망설이지 말고 새 세션을 여세요. 대화 기록을 조금 잃는 편이, 성능이 저하된 컨텍스트 속에서 계속 헤매는 것보다 훨씬 나아요.

\begin{tcbpractice}{직접 해보기: AI에게 한 번 “환자”가 되어 주기}
일부러 브라우저에서 작은 페이지 하나를 만들고 버그를 만들어 보세요(예: 버튼이 반응 없음). 그다음 이번 장의 “좋은 디버깅 설명” 형식대로 AI에게 한 번 완전하게 설명해 보세요: 원본 정보, 기대 vs 실제, 재현 단계, 확인 완료 항목. AI가 문제를 찾는 속도 차이를 직접 느껴 보면 “문제를 설명하는 것” 자체가 하나의 기술이라는 이유를 이해하게 될 거예요.

\end{tcbpractice}

\begin{tcbsummary}{이번 장 정리}
\begin{itemize}
  \item 프롬프트 3계층: 시스템 계층(기억 파일에 작성), 작업 계층(할 일+기술+제약), 피드백 계층(문제 지적+기대+맥락)
  \item 디버깅 4요점: 원본 정보 제공, 기대 vs 실제, 재현 단계, 확인 완료 항목
  \item 작은 단계 빠르게: 리스크 낮고 대화 횟수 많음, 신규 기능에 적합. 큰 단계 리팩터링: 리스크 높고 대화 횟수 적음, 테스트 뒷받침 필요
  \item 컨텍스트 관리 4가지: 핵심 정보 고정, 요약-이어가기, 한 대화에 한 작업, 과감하게 새 세션 열기
\end{itemize}

\end{tcbsummary}
